\section{Related Work}
\label{sec:related}
The proposed method connects three areas: learning from incomplete depth measurements, metric depth estimation with pretrained visual representations, and spatially selective refinement of an existing prediction. We distinguish these areas because SACG does not train a depth estimator from scratch. It refines an already strong, prompt-conditioned foundation model while keeping its backbone and decoder frozen.

\subsection{Depth Completion From Sparse Measurements}
The depth-completion task aims to reconstruct a dense metric field from sparse or low-resolution range measurements, often with guidance from a registered RGB image. Sparse-to-Dense~\cite{ma2018sparse} established a trainable encoder--decoder formulation for the joint processing of RGB and depth observations. The approach demonstrated that even relatively few metric samples can constrain dense inference when interpreted together with image structure. A key distinction from our setting is that a conventional completion network learns the entire depth mapping, whereas our adapter modifies only the residual of a pretrained metric estimator.

Convolutional spatial propagation networks (CSPN)~\cite{cheng2018cspn} treat dense depth refinement as repeated diffusion according to learned local affinities. Importantly, their propagation formulation can incorporate known sparse samples as constraints. NLSPN~\cite{park2020nlspn} extends propagation to non-local neighborhoods using learned confidence and affinities, showing the importance of spatially adaptive behavior when observed depth is unevenly distributed. These methods establish that learned spatial weighting and confidence-informed correction are well known in depth completion. Therefore, the novelty of SACG cannot be the mere use of a gate or the notion of confidence-guided refinement.

More recently, CompletionFormer~\cite{zhang2023completionformer} combines convolutional operators and transformers to leverage local details and non-local interactions. Such architectures are designed to propagate sparse geometry over large regions. Direct performance comparisons between these systems and a PromptDA adapter require consistent training data, sensor sparsity, resolution, and evaluation protocol. We do not infer superiority over published completion benchmarks from the paired validation results reported in this manuscript.

\subsection{Depth Foundation Models and Metric Depth Prompting}
Strong visual representation learning has changed how depth prediction systems are built. DINOv2~\cite{oquab2024dinov2} provides general visual features, whereas Depth Anything~\cite{yang2024depthanything} and Depth Anything V2~\cite{yang2024depthanythingv2} demonstrate the benefits of large-scale depth pretraining. These methods improve single-image geometry estimation but may still require metric anchoring when absolute depth is needed.

Prompt Depth Anything (PromptDA)~\cite{lin2025promptda} integrates low-cost LiDAR measurements into the depth foundation model's decoder and produces metric depth at substantially higher resolution than the input prompt. Its pretrained representation already captures rich visual and metric information. This creates a distinct adaptation setting: instead of discarding the pretrained prediction, we can inspect whether a small learned correction improves local residual errors. Our work uses the pretrained PromptDA ViT-L inference pathway without modifying its prompt-fusion layers, visual backbone, or decoder.

\subsection{Residual Prediction, Gating, and Confidence}
Given a coarse prediction $D_0$, residual refinement constructs an estimated correction $R$ and forms $D_0+R$. A spatially gated variant uses $D_0+G\odot R$, where $G$ is a learned coefficient at each location. The second formulation offers an explicit spatial modulation mechanism, but a sufficiently expressive residual CNN can potentially learn similar behavior without an independently parameterized gate.

The proposed SACG architecture separates the two roles. A residual CNN uses RGB, coarse depth, and intermediate frozen DPT features, while the gate network receives a normalized RGB--depth gradient-agreement map, coarse depth, and a learned field constructed from the locations of valid LiDAR measurements. This differs from direct propagation: SACG neither iteratively propagates measured depths nor explicitly enforces equality at sensor-observed pixels. The gate is also not calibrated predictive uncertainty, since no uncertainty labels, distributional likelihood, or calibration objective are used.

The distinction between LiDAR sample values and their spatial support is important. In the implemented simple CNN baseline, resized numerical LiDAR observations and the validity mask are direct inputs. SACG's residual CNN does not receive the numerical depth prompt directly; the frozen PromptDA receives those values upstream, and the SACG gate sees only a validity-derived field. The two trained heads consequently differ in feature availability in addition to gating. Their performance comparison is useful as a system-level control, but alone cannot isolate the gate's causal benefit.

\subsection{Datasets and Evaluation Perspective}
ARKitScenes~\cite{baruch2021arkitscenes} provides mobile RGB-D sequences and high-quality depth supervision relevant to the proposed use of low-cost metric prompts. ScanNet++~\cite{yeshwanth2023scannetpp} combines high-fidelity indoor scene geometry with commodity-sensor image and depth captures. Results depend on the particular scene splits, sensor alignment, valid-depth conventions, and reference-depth construction. Accordingly, we retain common frame manifests for paired comparisons and distinguish local dataset exports from official published benchmark settings.

The precise research question studied here is whether consistency-conditioned spatial modulation offers a repeatable benefit for supervised downstream refinement of an otherwise frozen prompt-conditioned depth foundation model. The existing experiments support only a preliminary, metric-specific answer; a more controlled study is needed to separate gate design from feature inputs and training-objective choices.
